% Physical unit 022; accepted Session 02.
\GordanUnitStart{022}{printed page 17}
\GordanTopFolio{17}
et integrando, ut in \S\ 11, et notando functiones
\[
 {dl\over d\psi_1}\qquad {dl\over dl_0}\qquad {dl'\over d\psi_1}\qquad {dl'\over dl_0}\qquad (01)\quad (11)
\]
indicibus iisdem, quibus respondentes $\psi$:
\begingroup\scriptsize
\[
\begin{aligned}
U=\left\{\begin{aligned}
&(01)_2\left\{\alpha_1\left({dl\over d\psi_1}\right)_2+\beta_1\left({dl\over dl_0}\right)_2\right\}^2
 +(11)_2\left\{\alpha_1\left({dl\over d\psi_1}\right)_2+\beta_1\left({dl\over dl_0}\right)_2\right\}
 \left\{\alpha_1\left({dl'\over d\psi_1}\right)_2+\beta_1\left({dl'\over dl_0}\right)_2\right\}\\
&-(01)_1\left\{\alpha_1\left({dl\over d\psi_1}\right)_1+\beta_1\left({dl\over dl_0}\right)_1\right\}^2
 -(11)_1\left\{\alpha_1\left({dl\over d\psi_1}\right)_1+\beta_1\left({dl\over dl_0}\right)_1\right\}
 \left\{\alpha_1\left({dl'\over d\psi_1}\right)_1+\beta_1\left({dl'\over dl_0}\right)_1\right\}\\
&+(01)\left\{\alpha_2\left({dl\over d\psi_1}\right)+\beta_2\left({dl\over dl_0}\right)\right\}^2
 +(11)\left\{\alpha_2\left({dl\over d\psi_1}\right)+\beta_2\left({dl\over dl_0}\right)\right\}
 \left\{\alpha_2\left({dl'\over d\psi_1}\right)+\beta_2\left({dl'\over dl_0}\right)\right\}\\
&-(01)_2\left\{\alpha_2\left({dl\over d\psi_1}\right)_2+\beta_2\left({dl\over dl_0}\right)_2\right\}^2
 -(11)_2\left\{\alpha_2\left({dl\over d\psi_1}\right)_2+\beta_2\left({dl\over dl_0}\right)_2\right\}
 \left\{\alpha_2\left({dl'\over d\psi_1}\right)_2+\beta_2\left({dl'\over dl_0}\right)_2\right\}
\end{aligned}\right.
\end{aligned}
\]
\endgroup

\begin{center}\S\ 13.\end{center}

Qua in aequatione constantes
\[
 \alpha_1,\quad \beta_1,\quad \alpha_2,\quad \beta_2,
\]
alternis relationibus inter se ita junctae sunt, ut in limitibus
\[
 \psi=\psi_1 \quad\text{et}\quad \psi=\psi
\]
$\delta l$ evanescat et intra limites sumto
\[
 \psi=\psi_2,
\]
unum tantum valorem habeat.

Cui rei respondet:
\[
 \alpha_1\left({dl\over d\psi_1}\right)_1+\beta_1\left({dl\over dl_0}\right)_1=0
 \qquad
 \alpha_2\left({dl\over d\psi_1}\right)+\beta_2\left({dl\over dl_0}\right)=0
\]
\[
 \alpha_1\left({dl\over d\psi_1}\right)_2+\beta_1\left({dl\over dl_0}\right)_2
 =\alpha_2\left({dl\over d\psi_1}\right)_2+\beta_2\left({dl\over dl_0}\right)_2,
\]
vel posito:
\[
\begin{aligned}
\alpha_1&=\left({dl\over dl_0}\right)_1
\left\{\left({dl\over d\psi_1}\right)_2\left({dl\over dl_0}\right)-\left({dl\over d\psi_1}\right)\left({dl\over dl_0}\right)_2\right\}\omega,\\
\beta_1&=-\left({dl\over d\psi_1}\right)_1
\left\{\left({dl\over d\psi_1}\right)_2\left({dl\over dl_0}\right)-\left({dl\over d\psi_1}\right)\left({dl\over dl_0}\right)_2\right\}\omega,\\
\alpha_2&=-\left({dl\over dl_0}\right)
\left\{\left({dl\over d\psi_1}\right)_1\left({dl\over dl_0}\right)_2-\left({dl\over d\psi_1}\right)_2\left({dl\over dl_0}\right)_1\right\}\omega,\\
\beta_2&=\left({dl\over d\psi_1}\right)
\left\{\left({dl\over d\psi_1}\right)_1\left({dl\over dl_0}\right)_2-\left({dl\over d\psi_1}\right)_2\left({dl\over dl_0}\right)_1\right\}\omega.
\end{aligned}
\]
Quibus valoribus ad aequationem supra propositam relatis, fit:
\[
 U=(11)_2\left\{\alpha_1\left({dl\over d\psi_1}\right)_2+\beta_1\left({dl\over dl_0}\right)_2\right\}
\left\{(\alpha_1-\alpha_2)\left({dl''\over d\psi_1}\right)_2+(\beta_1-\beta_2)\left({dl''\over dl_0}\right)_2\right\}.
\]
